% Configuración editorial para el PDF del libro
\usepackage{xcolor}
\usepackage{fontspec}
\usepackage{microtype}
\usepackage{fvextra}
\usepackage{tcolorbox}
\tcbuselibrary{skins,breakable}
\usepackage{titlesec}
\usepackage{fancyhdr}
\usepackage{booktabs}
\usepackage{longtable}
\usepackage{caption}
\usepackage{enumitem}
\usepackage{needspace}
\usepackage{etoolbox}

% Tipografías incluidas normalmente con TeX Live/TinyTeX
% Evitan depender de fuentes instaladas en Windows.
\setmainfont{TeX Gyre Pagella}
\setsansfont{TeX Gyre Heros}
\setmonofont[Scale=0.86]{TeX Gyre Cursor}

% Paleta visual
\definecolor{AzulLibro}{HTML}{1F4E79}
\definecolor{TurquesaLibro}{HTML}{168A91}
\definecolor{FondoCodigo}{HTML}{F4F6F8}
\definecolor{BordeCodigo}{HTML}{8BA6BF}
\definecolor{TextoCodigo}{HTML}{202124}
\definecolor{GrisEncabezado}{HTML}{5A6673}

% Composición general
\setlength{\parindent}{0pt}
\setlength{\parskip}{5pt plus 1pt minus 1pt}
\setlist[itemize]{topsep=4pt,itemsep=2pt,parsep=0pt}
\setlist[enumerate]{topsep=4pt,itemsep=2pt,parsep=0pt}
\emergencystretch=2em
\widowpenalty=10000
\clubpenalty=10000
\displaywidowpenalty=10000

% Títulos
\titleformat{\chapter}[display]
  {\normalfont\Huge\bfseries\color{AzulLibro}}
  {\chaptertitlename\ \thechapter}{10pt}{\Huge}
\titlespacing*{\chapter}{0pt}{-15pt}{25pt}
\titleformat{\section}
  {\needspace{4\baselineskip}\normalfont\Large\bfseries\color{AzulLibro}}
  {\thesection}{0.75em}{}
\titleformat{\subsection}
  {\needspace{3\baselineskip}\normalfont\large\bfseries\color{TurquesaLibro}}
  {\thesubsection}{0.75em}{}

% Encabezados y pies
\pagestyle{fancy}
\fancyhf{}
\fancyhead[LE]{\small\color{GrisEncabezado}\nouppercase{\leftmark}}
\fancyhead[RO]{\small\color{GrisEncabezado}\nouppercase{\rightmark}}
\fancyfoot[C]{\thepage}
\renewcommand{\headrulewidth}{0.35pt}
\renewcommand{\footrulewidth}{0pt}
\setlength{\headheight}{15pt}

% Bloques de código generados por Pandoc/Quarto
\DefineVerbatimEnvironment{Highlighting}{Verbatim}{
  commandchars=\\\{\},
  breaklines=true,
  breakanywhere=true,
  fontsize=\footnotesize,
  baselinestretch=1.03,
  frame=single,
  framesep=3mm,
  rulecolor=\color{BordeCodigo},
  fillcolor=\color{FondoCodigo},
  formatcom=\color{TextoCodigo},
  xleftmargin=1mm,
  xrightmargin=1mm
}

% Código sin resaltado
\fvset{
  breaklines=true,
  breakanywhere=true,
  fontsize=\footnotesize,
  baselinestretch=1.03
}

% Tablas y figuras
\renewcommand{\arraystretch}{1.15}
\captionsetup{font=small,labelfont=bf,justification=centering}
\AtBeginEnvironment{longtable}{\small}

% Nombres en español
\renewcommand{\contentsname}{Índice general}
\renewcommand{\chaptername}{Capítulo}
